\section{Discussion and Threats to Validity}
\label{sec:discussion}

\subsection{Three difficulties, measured}
\label{sec:discussion:difficulties}

\textbf{Root cause beyond symptoms.} This is the hardest part of the
task, and the benchmark shows it in a form that can be counted:
in \nHardPct{} of real reports no file the fix changed is named
anywhere in the text, and a method that matches the report against the
code it mentions scores zero there by construction. On those
\nEvalHard{} evaluation instances the retrieval layer reaches
\retrHardFive{} Acc@5 and \sys{} \hardFive{}. Two things recover a file
no log line names: incident memory, statistically, since a past fix
points at code the report never mentions; and reading the code,
causally, since the agent judges whether a file creates the bad state
or only reports it and walks back to the callers. The walk-back is one hop, not data-flow analysis. At \hardFive{}
the hard stratum is where future work should be judged, and it is not
solved.

\textbf{Noisy and incomplete logs.} Here the benchmark is kinder than
production: its logs are the excerpts reporters chose to paste, so a
human has filtered the noise once, and 30\% of reports carry a stack
trace. The pipeline groups repeated lines into events, ranks and labels
them, and finds the code that printed each by matching it against the
snapshot's logging statements; the agent cites those event ids. We did
not evaluate on synthetic noise or raw runtime bundles, which would
need segmenting by incident first.

\textbf{Evidence-to-code semantic retrieval.} Our results say where
log-and-code semantics have to sit. A small dense
retriever helps in the memory slot, where the task is comparing one
report with another, and loses to full-text BM25 in the code slot on
every one of the \nProjects{} projects, because matching a log line to
code is mostly identifier matching (\cref{sec:results:cost}). What
supplies the semantics in our system is a model that reads the code,
paired with a procedure that puts the right code in front of it.
The lower block of \cref{tab:all} is a ladder of how much code the same
model sees: path names only, 52.7\% Acc@1; paths with function names,
60.8\%; the files themselves, read whole or in whole functions, 67.1\%.

\subsection{What the explanation output is}
\label{sec:discussion:explanation}

The retrieval layer produces provenance rather than prose: for each
ranked file, which signals ranked it, the matching terms, and the past
incidents that voted for it. \sys{} produces an argument: for each file, the verdict and
the quoted line it rests on, and one chain from the reported symptom to
the named cause in which 57.3\% of steps cite one of the report's log
events, 85.7\% of them grounded, as an event id or a verbatim quote.
Both are checkable without being
trusted: a chain can cite a real line and still give the wrong cause,
so the evidence is a lead to verify, and no user study shows that
either form speeds engineers up.

\subsection{Threats to validity}
\label{sec:discussion:threats}

\textbf{Construct.} The answer key is the set of non-test files a fix
modified, a proxy for causation rather than a measurement of it. A
commit subject can name an issue it does not fix (reverts, merges and
long multi-commit efforts are removed); a fix can touch files
refactored in passing; and the repaired file is occasionally not where
the defect originated. With a
median of one gold file per instance the effect is bounded; accuracy
rises with fix size for every system because a larger fix offers more
files to hit, a loosening of the key rather than better localization. The module metric depends on a per-project depth
fixed before evaluation. Gold functions are derived from diff hunks and
miss a defect fixed by adding a function.

\textbf{Internal.} The prompts and step budgets of \sys{} (the
candidate table of thirty, the five files to shortlist, the eight to
examine, the 24{,}000-character file view) were developed on two
pilots drawn from the evaluation split, of 30 and 81 instances, both
disjoint from the \nSample-instance sample; no parameter was fitted to outcomes on evaluation instances, and
the pilot instances remain in the split (0.4\% of it). Scoring \sys{}
on the 26{,}429 answered instances outside the pilots changes nothing at the
paper's precision: file Acc@1 68.6\% and Acc@5 85.9\%, module 90.2\%,
region 53.8\% and 70.0\%. Two ablations
are missing. No run removes the logs from the reports, so the
contribution of the log lines themselves, as opposed to the rest of
the report, is not isolated; and incident memory is ablated in the
retrieval layer (\cref{tab:retrieval}) but not inside the agent, so
what the agent's \texttt{recall} evidence contributes on its own is not
measured. The agents run at temperature zero with a fixed seed, but
the model is not bit-reproducible under continuous batching.
On 20 instances run twice, the top-ranked file agreed on 18 on the same
machine and on 14 to 15 across two machines with identical software and
model checksums, and the module on all 20: the top file agrees across
repeats in 70\% to 90\% of cases, so differences of a point or two
between agent runs are inside the noise. Three seeds of \sys{} on
the \nSample-instance sample bound that noise in the aggregate: on the
998 instances all three answered, file Acc@1 spans 67.0\% to 67.9\%,
module 89.4\% to 90.2\% and region Acc@1 53.2\% to 53.3\%, at most
one point on every headline metric, while 55 to 61 instances change
their top file between two seeds. Four instances ended without an answer and are
scored as misses. Thinking was disabled for every model-based system, a
fair comparison that may understate the model.

\textbf{External.} Twenty-seven projects with disciplined issue-commit
linkage, mainly in four languages; a project without such linkage
cannot use the mining recipe directly. One model, so the result shows
what this architecture achieves with a 27B open-weights model, not how
it scales. Logs are
reporter-curated, as discussed above.

\textbf{Baselines.} LocAgent-style and CoSIL-style are our
reimplementations, not the published systems, and one check on the
ports goes the wrong way: CoSIL-style, whose design includes a step to
repair unparseable answers, left more answers unparseable than
LocAgent-style. Both ports leave 8\% to 10\% of instances without a
parseable answer, which counts as a miss; \cref{sec:results:agents}
reports the margins with those instances excluded as well. Pathidea,
the closest log-based prior work, was not run, not even on the Java
subset its path reconstruction is built for; that is a gap. Classic information-retrieval
localizers were not run; coverage-based localization needs execution
evidence absent at report time. The contamination probe
tests recitation from an issue key, not recognition when report and
code are shown, and the benchmark has no post-cutoff subset, so
memorisation cannot be excluded. The explanations were checked mechanically, for existence
and grounding; no human judged whether a chain's reasoning is correct,
and no practitioner has used the output.
